\documentclass[conference]{../../../../Conference-LaTeX-template_10-17-19/IEEEtran}
\usepackage{amsmath,amssymb,amsfonts,array,booktabs,tabularx,balance,url}
\newcolumntype{Y}{>{\raggedright\arraybackslash}X}
\newcommand{\R}{\mathbb{R}}

\title{What Is the Volume, Composition, and Significance of the Marine Deep Biosphere?\\
An Audit-First Research Program for Abundance, Activity, and Biogeochemical Inference}

\author{\IEEEauthorblockN{Research Proposal}
\IEEEauthorblockA{Evidence-bounded four-agent synthesis -- Environmental and astrobiology design only}}

\begin{document}
\maketitle

\begin{abstract}
The marine deep biosphere is often described as a vast, dark, cold, high-pressure microbial realm. That description is useful but insufficient for scientific accounting because it merges physically and biogeochemically distinct habitats: aphotic water, seafloor sediment and pore water, and upper oceanic crustal fluids. It also risks conflating four different quantities: habitat volume, standing cell abundance, functional genetic potential, and realized metabolic flux. This report surveys evidence for the extent, composition, and significance of these habitats and proposes DeepBio Census, a pressure-preserved, habitat-resolved inventory program. The evidence supports widespread microbial life in subseafloor sediment and crustal fluids; substantial bacterial and archaeal contributions that vary with energy and redox context; and relevance to carbon, sulfur, methane, nitrogen, mineral, and long-term ocean chemistry. It does not support a single universal volume, cell total, activity coefficient, climate effect, or extraterrestrial-life conclusion.

DeepBio Census integrates bathymetry and porosity, calibrated cell counts, multiomics, geochemistry, reaction-transport constraints, and contamination-aware provenance in a hierarchical model. It estimates separate posterior distributions for living-cell density, community composition, and redox rates within each habitat stratum before global aggregation. Mathematical design is included for inference and validation, but this is a proposal: no samples were collected, no genes were sequenced, no rate was measured, and no global census is claimed. The central contribution is an auditable route from local observations to bounded statements about the volume, composition, and significance of the marine deep biosphere.
\end{abstract}

\begin{IEEEkeywords}
deep biosphere, marine sediment, oceanic crust, deep ocean, Archaea, Bacteria, multiomics, biogeochemistry, hierarchical modeling, astrobiology
\end{IEEEkeywords}

\section{Introduction}
The ocean below the sunlit surface is one of Earth's largest connected environmental systems, but its living inventory cannot be captured by the phrase "deep ocean" alone. The term may refer to aphotic water masses, the seafloor, sediment thousands of meters below the seafloor, or fluids in basaltic crust. Each has different geometry, thermal history, pressure, permeability, residence time, carbon supply, and redox energy. A cell suspended in an abyssal water mass, a slow-growing cell in buried sediment, and a cell in a ridge-flank crustal fluid do not belong to a single homogeneous population merely because all experience darkness and hydrostatic pressure.

The question "What is the volume, composition, and significance of the deep biosphere?" links fundamental inventory science to climate, biodiversity, natural resources, and astrobiology. A defensible answer must separate what is measured. Volume is a geometric or pore-space quantity. Abundance is a count or biomass density integrated over specified spatial support. Composition describes taxa, genes, lipids, or genomes. Activity is a reaction rate. Significance is an interpretation of integrated activity in carbon, methane, sulfur, nitrogen, mineral, or habitability contexts. Each step needs a different observation model. A large population may have low per-cell turnover; a gene may be present but inactive; a local rate may not scale globally.

This report implements an evidence-bounded four-agent workflow. The Survey stage freezes literature that distinguishes sedimentary, water-column, and crustal habitats. The Idea stage selects DeepBio Census, a hierarchical inventory rather than a single-number extrapolation. The ExperimentDesign stage defines sampling strata, controls, measurements, models, and release rules. The Author stage organizes those inputs in IEEE form without adding a fictional field campaign, a global biomass estimate, or a claim of extraterrestrial detection.

The direct conclusion is that the marine deep biosphere is extensive, diverse, and biogeochemically consequential, but its exact global inventory is conditional on habitat definitions and measurement calibration. A near-term research target is not a more emphatic global total. It is a transparent posterior distribution that reports what is known, which compartment dominates uncertainty, and what new observation would reduce that uncertainty.

\section{Scope, Definitions, and Physical Context}
\subsection{A habitat-resolved definition}
For this proposal, the marine deep biosphere contains three linked habitat families:
\begin{enumerate}
\item \textit{Aphotic water column}: deep water masses, particles, and interfaces below persistent photosynthetic light;
\item \textit{Marine sediment and pore water}: seabed sediments from near-surface horizons through deeply buried strata, including methane- and sulfate-influenced zones;
\item \textit{Upper oceanic crust}: basalt-associated surfaces and fluids circulating through permeable crustal aquifers.
\end{enumerate}

The definition is broader than water deeper than a chosen isobath and narrower than all terrestrial subsurface life. It makes physical support explicit. For sediment, bulk volume, solids, pore water, and porosity must be reported separately. For crust, geometric rock volume is not the same as accessible fluid volume or biologically occupied surface area. For water, an oceanographic volume does not specify particle-associated versus free-living biomass.

The deep-sea conditions in the prompt require correction rather than rejection. Abyssal water near $4\,^\circ$C is common, but hydrothermal, subduction-zone, and sedimentary environments span wider temperatures. Hydrostatic pressure increases roughly one bar per ten meters of water depth, so 1,000 m corresponds to approximately 100 bar, while deeper settings experience substantially higher pressure. Pressure, temperature, darkness, and depth are useful covariates; they are not sufficient ecological coordinates. Energy supply, electron acceptors, diffusion, organic-carbon burial, and water-rock chemistry frequently provide the closer link to microbial function.

\subsection{Four estimands rather than one census}
Let $h$ index a habitat family, $s$ a spatial element, and $k$ a redox process. A transparent inventory distinguishes:
\begin{align}
V_h &= \int_{\Omega_h} 1\,dV, \label{eq:volume}\\
N_h &= \int_{\Omega_h} L_h(s)\,dV, \label{eq:cells}\\
F_{k,h} &= \int_{\Omega_h} R_{k,h}(s)\,dV, \label{eq:flux}
\end{align}
where $V_h$ is an explicitly defined accessible or occupied volume, $L_h(s)$ is living-cell density, and $R_{k,h}(s)$ is a realized volumetric rate. Community composition $P_h(s)$ is separate: it may be expressed through taxonomic relative abundance, genome-resolved pathway distribution, lipid signatures, or transcript profiles.

Equations \eqref{eq:volume}--\eqref{eq:flux} explain why a single number is rarely adequate. A microscopy count may contribute information about $L_h$. A bathymetric model may constrain $V_h$. A metagenome may inform $P_h$ and functional potential. A pore-water gradient or rate assay may constrain $R_{k,h}$. None identifies all other variables. In particular, $N_h$ is not $F_{k,h}$ multiplied by a universal coefficient. The unit mismatch is a central cause of overinterpretation in deep-biosphere summaries.

\subsection{Why depth-only extrapolation fails}
Depth covaries with pressure, temperature, oxygen, sediment age, and organic-matter quality, but the covariation differs by setting. In a continental margin, rapid burial can supply organic matter and create shallow sulfate-methane transitions. In an oligotrophic gyre, oxidants can persist deeply in slowly accumulating sediment. In crustal fluids, water-rock reaction and circulation can supply redox disequilibria independent of a sedimentary organic-carbon gradient. Inagaki \textit{et al.} reported contrasting microbial communities in methane-hydrate-bearing and hydrate-free deep marine sediments \cite{inagaki2006}; Hoshino \textit{et al.} found oxygen presence and organic-carbon concentration to be major factors in global marine-sediment composition \cite{hoshino2020}. These observations favor mechanistic habitat stratification over a one-dimensional depth model.

\section{Survey of Extent and Composition}
\subsection{Subseafloor sediment: large extent with heterogeneous occupancy}
Marine sediment covers most of Earth's surface, but the microbial census of sediment cannot be reduced to surface area times a uniform depth profile. Cell abundance and diversity commonly decline with burial depth and age, while local lithology, organic-carbon supply, methane, sulfate, temperature, and sedimentation create departures from a simple decline. Hoshino \textit{et al.} surveyed globally distributed cores and reported sediment depths from 0.1 to 678 m, finding that oxygen state and organic-carbon concentration were key determinants of taxonomic composition and diversity \cite{hoshino2020}. Inagaki \textit{et al.} documented microbial evidence in coal-bearing sediment approximately 1.5 to 2.5 km beneath the seafloor, with low indigenous cell concentrations below 1.5 km and pronounced lignite-layer contrasts \cite{inagaki2015}.

Archaea are an important component of the sedimentary biosphere, but an archaeal fraction is not a universal constant. Using microfluidic digital PCR across 221 core samples, Hoshino and Inagaki estimated that archaeal cells constituted 37.3 percent of total microbial cells in their global subseafloor-sediment dataset and inferred a global archaeal-cell count from that dataset \cite{hoshino2018}. The number rejects an older picture of bacteria alone dominating every sediment horizon. Its use is bounded: it is a model-based sediment estimate, not a fraction for deep seawater, crustal fluids, every individual core, or every physiological state.

Composition carries a time dimension. The transition from seabed to deep sediment can reflect selection from shallow community members rather than independent colonization at great burial depth. DNA can persist outside cells or adsorbed to mineral matrices, so an overall DNA pool is not equivalent to the living local community \cite{torti2015}. DeepBio Census treats extracellular DNA and extraction blanks as explicit error channels rather than a footnote after sequencing.

\subsection{Energy limitation and metabolic diversity}
The deep biosphere is frequently energy limited rather than metabolically inert. Lever \textit{et al.} synthesize laboratory and field evidence showing that microorganisms can persist over long periods under extremely low energy input, while cautioning against direct extrapolation from laboratory starvation to natural ecosystems \cite{lever2015}. Low cell-specific energy turnover may still integrate over enormous spatial extent and geological time, but rate must be measured or constrained rather than guessed from cell abundance.

In anoxic sediment, sulfate reduction is a dominant terminal pathway of organic-matter mineralization, and sulfur transformations interact with carbon, methane, iron, manganese, and nitrogen cycles \cite{wasmund2017}. Methane production, oxidation, and transport connect deeply buried processes to the seabed and overlying ocean, but pathways and rates vary by sedimentary zone. Global diffusive methane-flux work demonstrates the importance of domain-wide transport accounting \cite{egger2018}; it does not license attribution of every methane flux to a particular microbial population without site-specific evidence.

D'Hondt \textit{et al.} emphasize that subseafloor microbial activities affect major aspects of ocean chemistry and Earth's surface oxidation, while identifying diffusion rates and reaction affinities as primary rate controls \cite{dhondt2019}. The implication is direct: a deep-biosphere census should estimate energy and reaction environment together with cells and genes. The proposal treats a gene, transcript, lipid, geochemical gradient, and rate measurement as complementary evidence rather than interchangeable proxies.

\subsection{Upper oceanic crust and fluid circulation}
The upper oceanic crust is a large hydrogeologic reservoir and a difficult microbial habitat to sample without disturbance. Marine deep-biosphere reviews include sediments, pore water, basaltic crust, and fluids circulating through it because temperature, pH, donors, acceptors, and transport pathways change over millimeter-to-kilometer scales \cite{orcutt2013}. Crustal-fluid observations demonstrate microbial communities distinct from adjacent seawater and sediment, but the number of accessible sites remains sparse.

Time-series metagenomic and metatranscriptomic work in cold, oxic ridge-flank crustal fluids has revealed metabolically flexible communities with potential for organic-carbon oxidation, autotrophic pathways, and alternative electron acceptors such as nitrate and thiosulfate \cite{seyler2020}. This shows why a single taxonomic census is incomplete. Community composition can vary between horizons, and activity inferences require matching fluid chemistry, temporal context, and controls for drilling disturbance. The study supports long-term observatories, but not an unqualified global flux for the crustal biosphere.

\subsection{Astrobiology relevance}
Deep biosphere research is relevant to astrobiology because it tests how life persists with little light, limited energy, high pressure, water-rock chemistry, and long residence times. Colman \textit{et al.} describe deep, hot biosphere research as informative for Earth's evolution and the potential for life elsewhere, with hydrogen, methane, and carbon featuring prominently in many inferred metabolisms \cite{colman2017}. The valid transfer is conditional: if another world has liquid solvent, accessible redox disequilibria, suitable temperature-pressure conditions, and preservation pathways, Earth deep-biosphere observations inform a habitability hypothesis. They do not detect life on that world, identify a unique biosignature, or prove that an Earth lineage could survive there.

\section{Research Gaps and the Need for DeepBio Census}
\subsection{The volume-to-flux inference gap}
The first gap is geometric. A census must define which volume is biologically relevant: ocean water, sediment bulk, pore water, permeable crustal voids, or mineral surface. For a sediment element with bulk volume $dV_b$ and porosity $\phi(s)$, pore-water volume is $\phi(s)dV_b$, while solids occupy $(1-\phi(s))dV_b$. A cell density reported per cubic centimeter of wet sediment cannot be integrated as if it were per cubic centimeter of pore water. The same problem arises when crustal fluid volume is inferred from permeability and circulation rather than measured directly.

The second gap is calibration. Microscopy, flow cytometry, digital PCR, lipids, metagenomic coverage, and marker genes have different extraction efficiencies and different sensitivity to dead cells, extracellular DNA, and low-biomass contamination. A robust inventory must model these differences. A reported global cell total without sample-specific conversion assumptions may be numerically precise but scientifically underidentified.

The third gap is functional interpretation. Metagenomes reveal potential pathways; transcripts can indicate response or expression; neither is a direct flux meter. A sulfur-reduction gene supports the possibility of sulfur metabolism. A sulfate gradient and matched rate constraint make a stronger claim. A final carbon-cycle contribution requires integration over the spatial and temporal domain with uncertainty. DeepBio Census is designed around this ladder of evidence.

\subsection{Preservation, sparsity, and transfer gaps}
The fourth gap is signal preservation. Drilling fluids, reagent background, depressurization, storage, and incomplete recovery can each move the observed state away from the in situ state. This is especially consequential in low-biomass, low-energy settings, where a small introduced signal can dominate a molecular library. The fifth gap is coverage: most cores, observatories, and chemical records are concentrated in accessible settings, whereas global scaling needs representation across province, productivity regime, sediment age, redox state, water depth, and crustal circulation setting. The sixth gap is transfer. Earth deep-biosphere data are valuable analog evidence for icy-world and martian habitability questions, but a biological inference outside Earth requires independent measurements appropriate to that target.

\section{DeepBio Census: Selected Research Direction}
\subsection{Research question and hypothesis}
The selected Idea Agent direction is \textit{DeepBio Census}: a pressure-preserved, habitat-resolved Bayesian inventory of extent, composition, activity, and biogeochemical significance. The project asks whether geometry, pore-space accessibility, energy-redox environment, and measurement calibration predict abundance and activity better than depth alone. Its central hypothesis is:

\begin{quote}
\textit{Across aphotic water, marine sediment, and upper oceanic crust, energy-redox habitat stratification plus cross-platform calibration will explain observed living-cell density and process-rate variation more faithfully than a depth-only extrapolation.}
\end{quote}

This is an assertive but falsifiable direction. The hypothesis is not that every molecular observation will align, nor that one survey will close the global inventory. It predicts measurable improvement in out-of-sample predictive accuracy and posterior calibration when energy supply, electron donors and acceptors, organic-carbon context, porosity, temperature, and transport are represented explicitly. It also predicts that unresolved disagreement between observation channels will identify a methodological limitation or an ecological boundary rather than be averaged away.

\subsection{Causal and statistical structure}
For spatial element $s$ in habitat stratum $h$, let $X_h(s)$ contain measured geometry, porosity, temperature, pressure, organic-carbon context, dissolved species, mineralogy, transport, and redox variables. The latent quantities are living-cell density $L_h(s)$, composition vector $P_h(s)$, and realized process-rate vector $R_h(s)$. The inventory model is
\begin{align}
\log L_h(s) &= \alpha_h + X_h(s)^\top\beta_h + u_h(s), \label{eq:abundance-model}\\
\text{logit}\,P_{h,j}(s) &= \gamma_{h,j} + X_h(s)^\top\delta_{h,j} + v_{h,j}(s), \label{eq:composition-model}\\
\log R_{k,h}(s) &= \eta_{k,h} + Z_{k,h}(s)^\top\theta_{k,h} \nonumber\\
&\quad + \lambda_{k,h}\log L_h(s) + w_{k,h}(s), \label{eq:rate-model}
\end{align}
where $u$, $v$, and $w$ are spatially structured residuals and $Z$ contains process-specific chemical constraints. Equation \eqref{eq:abundance-model} represents abundance rather than an assumed biomass conversion; \eqref{eq:composition-model} represents a declared compositional assay; and \eqref{eq:rate-model} permits abundance to inform but not determine activity. The models are fit separately by habitat family before a global aggregation, avoiding an artificial exchangeability assumption among water-column, sediment, and crustal samples.

Observed data are linked to their latent targets through assay-specific observation models. For example, a calibrated microscopic count $Y^{(c)}$ may obey
\begin{equation}
Y^{(c)}_{h}(s) \sim \mathrm{NegBin}\!\left(q^{(c)}_{h}(s)L_h(s)+b^{(c)}_{h}(s),\,\kappa_c\right),
\label{eq:count-observation}
\end{equation}
where $q^{(c)}$ is recovery efficiency, $b^{(c)}$ is a blank- and contamination-informed background term, and $\kappa_c$ captures overdispersion. Analogous models are specified for digital PCR, lipid biomarkers, and sequencing readouts. This explicitly prevents a read count, a gene copy, and a living-cell count from being treated as the same unit.

The integrated posterior quantities are reported by stratum and globally:
\begin{equation}
N^{\mathrm{tot}}=\sum_h\int_{\Omega_h}L_h(s)\,dV,\qquad
F_k^{\mathrm{tot}}=\sum_h\int_{\Omega_h}R_{k,h}(s)\,dV.
\label{eq:global-integration}
\end{equation}
Equation \eqref{eq:global-integration} is a reporting rule, not a promise of a single precise total. Every estimate retains the support definition, uncertainty interval, source strata, and dominant extrapolation assumptions.

\subsection{Research objectives and decision rules}
The program has four objectives. First, establish a compatible physical inventory for water, sediment, and crustal fluid habitats. Second, quantify agreement and disagreement among abundance-assay channels using shared aliquots and controls. Third, distinguish functional potential from rate evidence by integrating multiomics with geochemistry and reaction-transport constraints. Fourth, provide an uncertainty-ranked global synthesis that identifies the most valuable next observation.

The primary model comparison is pre-registered conceptually: a depth-only baseline is compared with the energy-redox model in held-out habitat strata. DeepBio Census supports its central hypothesis when the expanded model improves predictive log score and calibration without systematic residual bias, and when its inferred rate ranges remain compatible with independent geochemical constraints. It is rejected or revised if depth-only performance is equivalent or better, if cross-platform abundance measurements cannot be reconciled within documented biases, or if inferred metabolic activity contradicts concentration gradients, mass balance, or direct rate constraints.

\section{Experiment Design}
\subsection{Study domain and sampling frame}
The research design uses a stratified, multi-province sampling frame rather than convenience cores alone. Each sample is classified by habitat family, depth or burial horizon, oxygen state, organic-carbon regime, temperature, pressure regime, sedimentation or circulation context, and major redox transition. Sediment strata include oxic oligotrophic sediment, organic-rich anoxic margin sediment, sulfate-methane-transition settings, and thermally elevated deep sediment. Crustal strata include cold oxic ridge-flank fluid, altered basalt-associated fluid, and hydrothermal-influenced end members. Aphotic-water strata are retained as a distinct family, not pooled with subseafloor material.

The design responds directly to evidence that oxygen and organic carbon organize sediment communities \cite{hoshino2020}, that temperature establishes meaningful subseafloor limits and zones \cite{heuer2020}, and that long-term deep-ocean observing infrastructure is needed to resolve variability and change \cite{levin2019}. Stratification is therefore based on mechanisms likely to affect survival and flux, while geographic province is included to expose rather than hide spatial transfer error.

\begin{table*}[t]
\caption{DeepBio Census measurement and validation plan. Each channel is assigned a target estimand and a boundary on interpretation.}
\label{tab:measurement-plan}
\centering
\footnotesize
\begin{tabularx}{\textwidth}{@{}p{0.16\textwidth}Y Y Y@{}}
\toprule
\textbf{Module} & \textbf{Primary measurements} & \textbf{Inference contribution} & \textbf{Required controls and validation} \\
\midrule
Physical habitat & Bathymetry, core interval, porosity, grain density, permeability, temperature, pressure, water depth, fluid-flow context & Defines bulk, pore-water, or accessible-fluid support for $V_h$ and the integration domain & Calibrated sensors; replicate porosity measurements; explicit wet/dry and bulk/pore-volume units; sample-registration record \\
Abundance calibration & Microscopy or flow counts, digital PCR, lipid biomarkers, biomass normalization & Constrains $L_h(s)$ and assay recovery terms in \eqref{eq:count-observation} & Field, procedural, drilling, extraction, and reagent blanks; matched aliquots; spike/recovery tests; independent laboratories where feasible \\
Composition and potential & 16S markers, shotgun metagenomes, genome-resolved pathways, extracellular-DNA fraction & Estimates declared taxonomic and functional-potential components of $P_h(s)$ & Negative controls, mock communities, contamination screening, read-depth reporting, separate extracellular-DNA treatment \\
Activity and chemistry & Dissolved oxygen, sulfate, methane, dissolved inorganic carbon, nitrate, metals, metabolites, rate constraints, reaction-transport parameters & Constrains $R_{k,h}(s)$ and tests whether potential is consistent with material balances & Duplicates, calibration standards, blank correction, redox-preserving handling, temporal repeats where systems vary \\
Integration and prediction & Hierarchical posterior model, spatial cross-validation, sensitivity analysis, mass-balance checks & Produces stratum-specific $N_h$, $F_{k,h}$, uncertainty, and information value & Hold out provinces and depths; compare against depth-only baseline; posterior-predictive checks; publish provenance and exclusions \\
\bottomrule
\end{tabularx}
\end{table*}

\subsection{Collection, preservation, and contamination control}
Sampling is pressure- and temperature-aware where the platform permits, with a continuous chain-of-custody record from retrieval to subsampling. For every physical sample, the registration record contains time, location, depth, pressure history, temperature, drilling or coring method, drilling-fluid tracers where applicable, storage interval, aliquot map, and all analytical transformations. The preferred comparison is not between unrelated cores analyzed by unrelated assays, but between matched aliquots taken from the same registered material.

Contamination control is a scientific variable. Field blanks, procedural blanks, reagent blanks, drilling-fluid controls, and extraction controls are processed through the same analytical steps as the primary samples. In low-biomass libraries, taxa and sequences associated with controls are flagged before biological interpretation. This does not imply that every repeated taxon is contamination; it requires a documented source comparison and a conservative decision rule. Extracellular DNA is quantified or separately treated because it can persist in sediments and otherwise inflate an inference about living cells \cite{torti2015}.

Native-state preservation also matters for chemistry. Degassing, oxygen intrusion, temperature change, and delay can alter dissolved gases, redox species, or expression signatures. The design therefore preserves split aliquots for cell counts, nucleic acids, lipids, and geochemistry, and it gives a sample reduced evidentiary weight when pressure, temperature, or storage deviation exceeds its declared tolerance. The result is not a claim that every in situ state can be perfectly recovered; it is a transparent record of how recovery uncertainty enters the analysis.

\subsection{Analysis pathway from counts to rates}
The first analytical layer estimates abundance using jointly modeled counts and independent calibration observations. A specimen is not promoted to a global cell contribution until its unit basis, recoveries, and blank-adjusted uncertainty are registered. The second layer estimates composition using taxonomic, genomic, and lipid assays whose target is named explicitly. Archaeal contributions are reported separately by habitat and stratum. The 37.3 percent archaeal fraction reported for a global subseafloor-sediment dataset \cite{hoshino2018} is used as context and a possible prior-sensitivity benchmark, never as a forced constant outside that dataset.

The third layer tests functional claims against chemical constraints. Candidate pathways are mapped to donors, acceptors, products, and stoichiometric consequences. For sulfur cycling, the presence of genes or transcripts is supplemented with sulfate, sulfide, and coupled-carbon evidence, consistent with the networked sulfur transformations described by Wasmund \textit{et al.} \cite{wasmund2017}. For methane, transport and concentration gradients are separated from biological source attribution, consistent with the domain-scale accounting problem examined by Egger \textit{et al.} \cite{egger2018}. Reaction-transport models impose conservation checks. A pathway may be retained as ``potential supported by molecular evidence'' even when the rate evidence is insufficient; it is not escalated to ``realized flux'' until the chemistry and model constraints agree.

Sensitivity analysis examines the effect of porosity maps, contamination priors, assay recovery, compositional reference databases, spatial covariance, and the choice of global extrapolation weights. Cross-validation withholds entire provinces, horizons, and redox classes in turn. This is harder than randomly withholding individual aliquots, but it tests the actual question: whether a census generalizes to environments not yet directly sampled.

\subsection{Ethical, operational, and safety boundary}
This design is an environmental research plan. It does not authorize autonomous drilling, collection, incubation, sequencing, deep-sea deployment, or intervention in protected marine areas. Cruise planning, permitting, sample handling, biosafety, shipping, site access, community consultation, and environmental-impact assessment require qualified human institutions and applicable jurisdictional review. Data release should protect sensitive locations where disclosure could increase extraction pressure or compromise protected habitats. The computational inference can be executed only after human approval of registered inputs, provenance, and protocol compliance.

\section{Expected Outcome Branches and Significance}
\subsection{Evidence-conditioned conclusions}
The study is designed to support a decision tree rather than a prewritten result. If the energy-redox model outperforms depth-only extrapolation and remains stable under calibration sensitivity analyses, the result supports a habitat-resolved explanation of abundance and activity. If its advantage occurs in sediment but not in crustal-fluid strata, the correct conclusion is domain-specific: the model has explanatory value in one habitat family and requires revision in another. If cross-platform measurements disagree beyond their modeled biases, the product is a measurement-gap diagnosis, not a blended global estimate.

For functional composition, concordant genome-resolved pathways, transcripts or lipids, and geochemical constraints support a bounded inference about an active process in the sampled stratum. Molecular evidence without a rate constraint supports potential only. Conversely, a chemical gradient without resolved biological identity may support a process-level rate estimate while leaving the responsible community ambiguous. This separation keeps the report useful in the common case where different evidence layers have unequal resolution.

\subsection{Scientific and societal significance}
The deep biosphere matters because slow reactions distributed over large areas and long times can influence oceanic chemical inventories. Subsurface life has been linked to major biogeochemical effects through the interaction of transport, reaction affinity, and energy supply \cite{dhondt2019}. By attaching uncertainty to volume, abundance, and rate separately, DeepBio Census can clarify which claims are globally robust and which depend on sparsely sampled compartments. This is more actionable for carbon- and methane-cycle assessment than a single headline cell total.

The program also improves biodiversity accounting. A habitat-resolved record can distinguish sediment-associated lineages, free-living fluid communities, and particle-associated deep-water populations. It enables reanalysis when taxonomy changes because raw sample provenance, controlled assays, and model assumptions remain linked. For natural-resource and conservation decisions, the same record identifies regions where physical disturbance would intersect with poorly characterized but chemically unusual microbial systems.

For astrobiology, the project contributes an experimentally disciplined analogue framework. Deep-ocean and subseafloor environments demonstrate how biological hypotheses should be tested under limited light and constrained energy, but the transfer is conditional on environmental comparability. The output for planetary science is a set of measurable habitability and preservation criteria, not a claim that an ocean world or planet contains life. This boundary makes the analogue stronger: it identifies what a future mission must measure to discriminate a biological hypothesis from an abiotic one.

\section{Risks, Limitations, and Human Review Requirements}
The principal scientific risk is underidentification: multiple assay biases or ecological processes can explain the same observation. The mitigation is not to claim certainty, but to report the competing explanations, maintain assay-specific observation models, and collect the next sample where the explanations diverge most. Sparse spatial coverage may produce overconfident global posteriors. Stratified validation and province-level withholding are mandatory before global integration. Poor preservation and low biomass can induce false molecular signals; the mitigation is matched controls, documented processing order, and exclusions exposed in the provenance record.

Some limitations are irreducible in this proposal. A core or fluid sample is a time-stamped perturbation of a dynamic system. Rate incubations may not represent geological time scales. Taxonomic databases are incomplete, and function annotation can be ambiguous. A model may accurately predict sampled settings yet fail at unobserved hydrogeologic regimes. These limitations define the language of the conclusions: conditional credible intervals and scope-specific claims, not universal constants.

Human review is required before interpreting high-consequence results, releasing sensitive geospatial data, accepting field protocols, or translating any analogue result to a planetary-life statement. Experts in marine microbiology, geochemistry, physical oceanography, statistics, environmental permitting, and data governance should review the relevant stage. The proposal contains no observed results and makes no claim that an intervention was performed.

\section{Conclusion}
The volume, composition, and significance of the marine deep biosphere cannot be answered by multiplying a global water volume by a representative cell density or by treating DNA as an activity meter. The evidence supports a broader and more useful conclusion: deep sediments, aphotic waters, and oceanic crust host microbial systems whose abundance and functions are structured by physical support, energy availability, redox chemistry, and transport. Their exact global inventory remains uncertain because those drivers and the measurements themselves vary across habitats.

DeepBio Census converts that uncertainty into a testable research program. It defines habitat volume before integration, calibrates multiple abundance channels, separates functional potential from realized rate, and requires posterior predictions to survive spatially meaningful validation. Its contribution is therefore not an invented final count. It is an auditable path to better counts, composition estimates, flux constraints, and astrobiology-relevant habitability criteria. The decisive next advance is a coordinated, pressure-aware, provenance-rich observing network that measures the deep biosphere as a set of linked but distinct environments.

\appendices
\section{Census Data and Provenance Schema}
Each observation is stored with a persistent sample identifier and links to a habitat definition, spatial coordinates or protected-location token, water depth, core or fluid interval, collection method, pressure-temperature history, time stamp, physical units, aliquot identifier, analytical method, controls, processing version, and quality flag. The registry stores a declared estimand for every value: physical support, living-cell abundance, taxonomic composition, functional potential, chemical concentration, or realized rate. It also stores the relevant numerator, denominator, unit conversion, detection limit, calibration result, and uncertainty distribution. This schema prevents a pore-water density from silently becoming a bulk-sediment density and permits later reinterpretation as assays or reference databases improve.

Model outputs are versioned with the exact input records, exclusion decisions, spatial weights, priors, diagnostic summaries, and integration domain. A result card states its habitat family, temporal window, evidence level, posterior interval, and known confounders. A release candidate is rejected when its evidence path cannot be reconstructed from the sample registry.

\section{Claim--Evidence Matrix}
\begin{table}[h]
\caption{Claim categories and minimum evidence required for release.}
\label{tab:claim-evidence}
\centering
\footnotesize
\begin{tabularx}{\columnwidth}{@{}p{0.24\columnwidth}Y@{}}
\toprule
\textbf{Claim category} & \textbf{Minimum evidence and prohibited shortcut} \\
\midrule
Habitat volume & Geometry plus a declared physical support; do not equate bulk rock with accessible fluid volume. \\
Living abundance & Calibrated count evidence with controls and stated unit; do not infer directly from sequence reads. \\
Functional potential & Taxonomic/genomic/lipid evidence with assay provenance; do not call a gene a measured rate. \\
Realized process rate & Chemical or direct-rate constraint compatible with mass balance; do not infer from cell abundance alone. \\
Global significance & Stratum-specific integration with spatial validation and uncertainty; do not generalize one core to the ocean. \\
Astrobiology relevance & Environmental analogue and explicit boundary; do not treat Earth data as extraterrestrial detection. \\
\bottomrule
\end{tabularx}
\end{table}

\section{Conditional Integration Template}
For each habitat family, the final synthesis reports $(V_h,N_h,F_{k,h})$ with a posterior interval, a physical definition of $V_h$, the observations informing $L_h$ and $R_{k,h}$, and the fraction of the integration domain represented by direct sampling. A decision dashboard ranks each stratum by contribution to total variance. If a small number of unsampled crustal or sedimentary strata dominate the uncertainty, the next expedition is allocated there rather than to already densely sampled stations. This makes the census adaptive while retaining a fixed evidence contract.

\begin{thebibliography}{99}
\bibitem{hoshino2018}
T. Hoshino and F. Inagaki, ``Abundance and distribution of Archaea in the subseafloor sedimentary biosphere,'' \textit{The ISME Journal}, vol. 13, pp. 227--231, 2019, doi: 10.1038/s41396-018-0253-3.

\bibitem{hoshino2020}
T. Hoshino \textit{et al.}, ``Global diversity of microbes from the marine sedimentary biosphere,'' \textit{Proceedings of the National Academy of Sciences}, vol. 117, no. 44, pp. 27587--27597, 2020, doi: 10.1073/pnas.1919139117.

\bibitem{inagaki2006}
F. Inagaki \textit{et al.}, ``Microbial community in a sediment-hosted methane hydrate reservoir: implications for subsurface life in methane-rich natural environments,'' \textit{Proceedings of the National Academy of Sciences}, vol. 103, no. 38, pp. 14189--14194, 2006, doi: 10.1073/pnas.0511033103.

\bibitem{inagaki2015}
F. Inagaki \textit{et al.}, ``Exploring deep microbial life in coal-bearing sediment down to $\sim$2.5 km below the ocean floor,'' \textit{Science}, vol. 349, no. 6246, pp. 420--424, 2015, doi: 10.1126/science.aaa6882.

\bibitem{orcutt2013}
B. N. Orcutt \textit{et al.}, ``Microbial activity in the marine deep biosphere: progress and prospects,'' \textit{Frontiers in Microbiology}, vol. 4, art. 189, 2013, doi: 10.3389/fmicb.2013.00189.

\bibitem{lever2015}
M. A. Lever \textit{et al.}, ``Life under extreme energy limitation: a synthesis of laboratory- and field-based investigations,'' \textit{FEMS Microbiology Reviews}, vol. 39, no. 5, pp. 688--728, 2015, doi: 10.1093/femsre/fuv020.

\bibitem{dhondt2019}
S. D'Hondt \textit{et al.}, ``Presence of oxygen and aerobic communities from seafloor to basement in deep-sea sediments,'' \textit{Nature Communications}, vol. 10, art. 2555, 2019, doi: 10.1038/s41467-019-11450-z.

\bibitem{seyler2020}
L. M. Seyler \textit{et al.}, ``Long-term development of the crustal microbiome in a subseafloor aquifer,'' \textit{The ISME Journal}, vol. 15, pp. 139--153, 2021, doi: 10.1038/s41396-020-00843-4.

\bibitem{heuer2020}
V. B. Heuer \textit{et al.}, ``Temperature limits to deep subseafloor life in the Nankai Trough subduction zone,'' \textit{Science}, vol. 370, no. 6521, pp. 1230--1234, 2020, doi: 10.1126/science.abd7934.

\bibitem{wasmund2017}
K. Wasmund, M. Mu{\ss}mann, and A. Loy, ``The life sulfuric: microbial ecology of sulfur cycling in marine sediments,'' \textit{Environmental Microbiology}, vol. 19, no. 2, pp. 487--509, 2017, doi: 10.1111/1758-2229.12538.

\bibitem{egger2018}
M. Egger \textit{et al.}, ``Global diffusive fluxes of methane in marine sediments,'' \textit{Nature Geoscience}, 2018, doi: 10.1038/s41561-018-0122-8.

\bibitem{torti2015}
A. Torti, M. A. Lever, and B. B. J\o{}rgensen, ``Origin, dynamics, and implications of extracellular DNA pools in marine sediments,'' \textit{Marine Genomics}, vol. 24, pp. 185--196, 2015, doi: 10.1016/j.margen.2015.08.007.

\bibitem{colman2017}
D. R. Colman \textit{et al.}, ``The deep, hot biosphere: twenty-five years of retrospection,'' \textit{Proceedings of the National Academy of Sciences}, vol. 114, no. 27, pp. 6895--6903, 2017, doi: 10.1073/pnas.1701266114.

\bibitem{levin2019}
L. A. Levin \textit{et al.}, ``Global observing needs in the deep ocean,'' \textit{Frontiers in Marine Science}, vol. 6, art. 241, 2019, doi: 10.3389/fmars.2019.00241.
\end{thebibliography}

\end{document}
